\section{Methods}

The systematic scoping-review protocol was frozen on 1 October 2026 and is specified in \texttt{review\_protocol.yaml}. This document is the reference for all subsequent eligibility decisions; changes must be logged and justified rather than inferred from automated screening scores.

\subsection{Eligibility}

Eligible reports must study humans using an interactive digital task, game, digital intervention, or instrumented computer-based paradigm; include an identifiable perturbation or challenge; and report or permit reconstruction of a temporally defined post-perturbation recovery process. Simple pre--post change, generic response time, or symptom improvement without an event-linked recovery definition is insufficient.

No diagnosis is required. ADHD is an important subgroup but not an inclusion requirement.

\subsection{Information sources}

The planned sources are MEDLINE/PubMed, PsycINFO, Scopus, Web of Science Core Collection, IEEE Xplore, and ACM Digital Library. Backward and forward citation searching will be recorded separately.

\subsection{Core search concepts}

Searches combine three concept blocks: digital/interactive task terms; perturbation or challenge terms; and recovery-dynamics terms such as recovery latency, return to baseline, stabilization, time-to-recovery, post-error recovery, or post-stressor recovery. Database-specific queries will be stored verbatim in the search log.

\subsection{Screening}

Deduplicated records will undergo title/abstract and full-text screening against prespecified criteria. Exclusion reasons at full text will be coded. Any protocol clarification made after screening starts will be recorded as a deviation.

\subsection{Preliminary electronic-search material and evidence integrity}

An externally prepared SciSpace export received on 28 October 2026 was audited as \emph{preliminary search material}, not as completed searching or screened evidence. The supplied combined file contains 577 records, with 26 automated candidates at the stage labelled full-text screening and 23 automatically selected records. The supplied screening notes state that full text was unavailable for all 23 selected records. The apparent selection counts must therefore not be entered as validated PRISMA screening or inclusion totals. The automated criteria also added a requirement for incremental predictive validity that is absent from the frozen inclusion rules; this is an extraction and appraisal topic, not an eligibility gate. Study-level reassessment, provenance, seed-search failures, and DOI corrections are recorded in the repository audit under \texttt{research/r01\_scispace\_2026-10-28/}. The planned database-specific searches and independent full-text assessments remain outstanding.

\subsection{Data extraction}

Extraction covers population, task/intervention, challenge definition, reference state, temporal origin, recovery endpoint, threshold, stability/persistence rule, censoring, missingness, within-person normalization, sampling rate, sensor/task modality, reliability/validity evidence, relation to external outcomes, preregistration, and availability of code/data.

\subsection{Synthesis}

The primary synthesis is a structured taxonomy rather than pooled efficacy. Frequencies and cross-tabulations will summarize measurement choices where meaningful. Conceptual distinctions will be retained when aggregation would erase materially different recovery definitions.

\subsection{Provisional study-flow visualization and reporting boundary}

For transparent audit of the supplied SciSpace export, an editable preliminary PRISMA-style flow diagram is stored in \texttt{figures/prisma\_flow\_preliminary.svg}, with its machine-readable counts, original CSV hashes and regeneration script documented under \texttt{research/r01\_scispace\_2026-10-28/prisma/}. The visualized counts are 577 composite source rows, 551 automated abstract exclusions, 26 automated candidates, three subsequent automated exclusions and 23 machine-retained records with no available PDF. These numbers describe tool labels, not independently validated database identification, reports assessed at full text or studies included. The latter formal PRISMA quantities remain explicitly unestablished. The figure is an audit artifact and is not yet the final PRISMA 2020 study-selection figure for the manuscript.
